% Adaptive Retrieval for Repository-Level Code Completion
% 8-10 page two-column paper skeleton (IMPLEMENTATION_GUIDE §16 Phase 8 outline).
%
% Compile with:   latexmk -pdf main.tex
% TODO comments mark placeholders for results that depend on Phase 6 runs.

\documentclass[10pt,twocolumn]{article}

\usepackage[margin=0.8in]{geometry}
\usepackage{times}
\usepackage{graphicx}
\usepackage{booktabs}
\usepackage{amsmath}
\usepackage{hyperref}
\usepackage{xcolor}
\usepackage{url}
\usepackage{microtype}
\usepackage{algorithm}
\usepackage{algorithmic}

\hypersetup{
    colorlinks=true,
    linkcolor=blue,
    citecolor=blue,
    urlcolor=blue,
}

\newcommand{\todo}[1]{\textcolor{red}{[TODO: #1]}}
\newcommand{\TRAG}{T_{\text{RAG}}}
\newcommand{\TACC}{T_{\text{ACC}}}

\title{Adaptive Retrieval for Repository-Level Code Completion:\\
A Cascade of Uncertainty and Static Analysis}

\author{Project Group 17\\
% Authors anonymised for double-blind review
}

\date{}

\begin{document}
\maketitle

\begin{abstract}
Retrieval-augmented generation (RAG) improves accuracy on repository-level
code completion but spends compute on every instance, including ones where
retrieval is unhelpful or harmful. CARD~\cite{zhang2024card} addresses this
with a learned uncertainty gate that skips retrieval when the model's
zero-shot prediction looks confident. We observe that CARD's uncertainty
signal misses a specific class of failures: \emph{confidently-generated
identifiers that don't actually exist in the repository}. We add a parallel
gate based on static analysis of the model's prediction, which fires
retrieval on instances where any used identifier fails to resolve in the
file's scope, the repository's symbol table, or the language's builtins. The
combined cascade preserves CARD's retrieval savings while
\todo{XX\%} lowering the hallucination rate on CrossCodeEval-Python
(McNemar's $p < $ \todo{XX}), and matching always-retrieve on edit
similarity at \todo{XX\%} of the retrieval cost.
\end{abstract}

\section{Introduction}

Code-completion models trained on large code corpora can reliably continue
short snippets of code, but performance degrades sharply when the right
completion depends on context elsewhere in the same repository --- a
function defined two files away, a class method whose signature lives in a
sibling module, an internal API whose name is conventional only to that
project. Retrieval-augmented generation closes much of this gap by
prepending relevant snippets to the prompt, but it does so indiscriminately,
spending compute on instances where the in-file context already suffices or
where retrieved snippets actively distract the model.

\textbf{Adaptive retrieval} addresses this by deciding, per instance,
whether to retrieve. CARD~\cite{zhang2024card} learns a regressor over
per-token log-probabilities and entropies to predict the zero-shot
output's edit similarity, and triggers retrieval when this prediction is
below a threshold. Our work starts from the observation that CARD's
signal --- aggregated logit confidence --- is blind to a specific failure
mode that is both common and pernicious: hallucinated identifiers.

A model can generate \texttt{self.client.send(\textcolor{red}{msg})} with
high token-level probability across the entire span, even when
\texttt{msg} is not bound anywhere in the surrounding code. CARD will see
high confidence and skip retrieval; the prediction will be silently
wrong; the user will see a plausible-looking line of code that fails at
import time.

\textbf{Contributions.} We add a static-analysis cascade stage to CARD's
decision flow. The new stage runs only when CARD says ``skip'' and triggers
retrieval on any of three orthogonal hallucination classes:
\textbf{(i)} an \emph{unresolved} reference --- an identifier in the
prediction's AST that resolves nowhere in scope, the repository, or the
language's builtins; \textbf{(ii)} a \emph{signature mismatch} --- a call
to a known repo function with the wrong arity or an unknown keyword
argument; and \textbf{(iii)} a \emph{wrong-import} attribution --- a
\texttt{from X import Y} statement where Y actually lives in module Z and
is not re-exported by X. A weaker fourth signal, \emph{cross-file}
references, is studied as an ablation: a cross-file identifier in the
prediction is evidence the model already recovered the right name, so
firing retrieval there spends budget on a case the parametric memory
already nailed and risks the retrieved-context distractor effect
documented in CARD~\cite{zhang2024card} and FLARE~\cite{jiang2023flare}.
The cascade is asymmetric by construction: static analysis can only
\emph{add} retrievals to CARD's decisions, never remove them, which bounds
the worst-case retrieval count to always-retrieve and frames the central
research question cleanly: \emph{does the extra retrieval budget reduce
hallucinations?}

We evaluate on CrossCodeEval-Python and RepoEval (line and API
subsets), comparing the cascade against vanilla CARD, always-retrieve,
never-retrieve, a static-only ablation, and an oracle upper bound. Our
empirical results show \todo{XX}.

\section{Background and related work}

\paragraph{Repository-level code completion.} \todo{1--2 paragraphs.
Reference RepoCoder~\cite{zhang2023repocoder} which introduces the iterative
retrieval baseline and RepoEval benchmark; CrossCodeEval~\cite{ding2023cce}
which curates instances explicitly requiring cross-file context;
Repoformer~\cite{wu2024repoformer} which proposes self-supervised retrieval
gating. Contrast the three approaches' assumptions: RepoCoder always
retrieves and iterates; Repoformer learns when not to retrieve via a
self-evaluation token; CARD learns the same with a separately-trained
regressor.}

\paragraph{Uncertainty for adaptive retrieval.} \todo{1 paragraph. CARD's
13-D feature vector and LightGBM regressor; FLARE~\cite{jiang2023flare}
which uses look-ahead probabilities for active retrieval in long-form
generation; the broader literature on calibration in generative models.}

\paragraph{Hallucination detection in code.} \todo{1 paragraph. Existing
work on code hallucinations (typically post-hoc, using exact-match or
execution); our position: static analysis of the prediction is a cheap,
language-agnostic signal that can be used \emph{pre-execution} to trigger
retrieval, sidestepping the need to actually run the predicted code.}

\section{Method}

\subsection{The cascade}

Given the in-file context $X = (X_{\text{left}}, X_{\text{right}})$ and the
repository $R$, the cascade decides the final prediction $\hat{y}$ in three
asymmetric stages:

\begin{algorithm}[H]
\caption{The cascade pipeline (single-RAG).}
\begin{algorithmic}[1]
\STATE $\hat{y}_0, \text{logits}_0 \gets G(X_{\text{left}}, X_{\text{right}})$
\STATE $\hat{s}_0 \gets \text{Estimator}(\text{features}(\text{logits}_0))$
\IF{$\hat{s}_0 < \TRAG$}
    \STATE $CC \gets \text{BM25}(X_{\text{left}}, R)$
    \STATE \textbf{return} $G(X_{\text{left}}, X_{\text{right}}, CC)$, \textbf{reason}=\texttt{card}
\ENDIF
\STATE $\text{sa} \gets \text{StaticAnalyzer}(\hat{y}_0, X_{\text{left}}, X_{\text{right}}, R)$
\IF{$\text{sa.fires}$}
    \STATE $CC \gets \text{BM25}(X_{\text{left}}, R)$
    \STATE \textbf{return} $G(X_{\text{left}}, X_{\text{right}}, CC)$, \textbf{reason}=\texttt{static\_\{unresolved$|$signature$|$import$|$crossfile\}}
\ENDIF
\STATE \textbf{return} $\hat{y}_0$, \textbf{reason}=\texttt{none}
\end{algorithmic}
\end{algorithm}

\subsection{CARD's uncertainty gate}

We reimplement CARD's single-RAG variant
faithfully~\cite{zhang2024card}: a LightGBM regressor predicts edit
similarity from a 13-dimensional feature vector built from per-token
probabilities and entropies (max, min, mean, std, product, geometric
mean for each, plus generation length). The Estimator is trained on
\todo{50k/250k} (X, y) pairs sampled from The Stack
subset~\cite{kocetkov2022stack}, with the same generator that will run
inference. \todo{Brief note on calibration MSE achieved vs the paper's 0.07.}

\subsection{Static analysis of the prediction}

The cascade's static stage decomposes ``does this prediction look like a
hallucination?'' into three independent signals over the prediction's
abstract syntax tree, plus a fourth (cross-file) retained only for the
ablation.

\paragraph{Unresolved identifiers.} For each identifier used in
$\hat{y}_0$ we classify the name as \textbf{builtin} (Python builtin or
visible at the hole), \textbf{in-file} (in scope at the completion
position), \textbf{cross-file} (defined somewhere in $R$ but not visible
in-file), or \textbf{unresolved} (nowhere). The analyzer uses tree-sitter
for fault-tolerant parsing and tracks scope through imports (including
aliased and \texttt{from X import Y} forms), function parameters,
for-loop and comprehension variables, walrus expressions, \texttt{with}
aliases, and lambda parameters. An \textbf{unresolved} name is a
near-certain hallucination: the model produced an identifier that exists
nowhere we can see.

\paragraph{Signature mismatches.} For each \texttt{call} node in
$\hat{y}_0$ whose callee can be resolved to a unique repository signature
(bare-identifier calls, module-qualified calls of the form
\texttt{mod.func(...)} when \texttt{mod} is an imported repo module, and
class-qualified method calls \texttt{Cls.m(...)}), we compare the call's
positional-argument count against the declared range
$[\text{min\_positional}, \text{max\_positional}]$ (deriving the bounds
from the signature's defaults and \texttt{*args}) and check every keyword
argument against the callee's parameter set (unless the callee declares
\texttt{**kwargs}). Calls whose callee can't be uniquely resolved
--- method calls on untyped receivers (\texttt{self.x}, \texttt{obj.x}),
multi-step attribute chains, calls whose callee is a call result,
external library callees, or callees decorated with signature-changing
wrappers --- are skipped. This conservative abstention preserves
precision over recall.

\paragraph{Wrong-import attributions.} For each
\texttt{from X import Y} statement in $\hat{y}_0$ we check that
\texttt{Y} is either defined in module \texttt{X} or is re-exported by
\texttt{X}'s \texttt{\_\_init\_\_.py} (via direct re-import or an
\texttt{\_\_all\_\_} declaration). When \texttt{Y} is defined elsewhere
in the repository as module \texttt{Z $\neq$ X} we flag the import as
\texttt{wrong\_origin}. A parallel check applies to module-qualified
attribute usage: for an in-file \texttt{import X as M} binding,
\texttt{M.Y} is flagged when \texttt{Y} lives outside module \texttt{X}.
Star imports and external modules trigger abstention.

\paragraph{Cross-file references (ablation only).} A cross-file name
in the prediction is evidence the model already recovered the right
identifier from parametric memory or in-file hints, so firing retrieval
there spends budget on a case the model already nailed and risks the
retrieved-context distractor effect. We therefore disable the
\texttt{fire\_on\_crossfile} flag by default but retain it as a
configurable knob so Section~\ref{sec:ablations} can quantify its
marginal contribution.

\paragraph{Trigger-reason precedence.} When multiple signals fire on the
same prediction, we report the strongest:
\texttt{static\_unresolved} $>$ \texttt{static\_signature} $>$
\texttt{static\_import} $>$ \texttt{static\_crossfile}. All diagnostic
payloads (unresolved names, call issues, import issues, cross-file
names) are still surfaced in the per-instance JSONL so the analysis in
Section~\ref{sec:analysis} can reconstruct any of the lower-priority
classes post-hoc.

\subsection{Implementation}

\todo{Brief paragraph: 7B-parameter code LMs (Qwen2.5-Coder or
CodeLlama-7B-hf), BM25 over 20-line stride-10 chunks, FIM prompt
formatting per family, generation cache keyed on (model, prompt, max\_tokens)
that shares the zero-shot pass across C1/C3/C4. Roughly XX lines of Python
across XX modules, released at \url{TODO}.}

\section{Experiments}

\subsection{Setup}

\paragraph{Datasets.} CrossCodeEval-Python (2{,}665 instances,
\cite{ding2023cce}) is the primary evaluation. RepoEval-line and -api
(1{,}600 instances each, \cite{zhang2023repocoder}) are used to validate
the CARD reimplementation against the paper's published numbers and as a
secondary evaluation surface.

\paragraph{Configurations.} We compare six configurations: (C1)
no-retrieve, (C2) always-retrieve with BM25 top-10, (C3) vanilla CARD,
(C4) our cascade, (C5) static-only ablation (retrieve iff static analysis
fires on the zero-shot output), and (C6) an oracle that runs both C1 and
C2 per instance and keeps whichever has higher edit similarity against the
ground truth.

\paragraph{Metrics.} Following~\cite{ding2023cce} and~\cite{zhang2024card}:
exact match (EM), edit similarity (ES), and identifier-F1 (IdF1) for
accuracy. For efficiency: percent retrieval and mean latency (ms).
For hallucination: \emph{repository-symbol precision} --- the fraction
of identifiers in the prediction that resolve somewhere in scope ---
and a per-instance \emph{hallucination flag} (true iff at least one
unresolved identifier).

\subsection{CARD reproduction}

\todo{Table comparing our CARD reimplementation against the paper's Table 3
on RepoEval-line and RepoEval-API. Target: within $\pm 1\%$ ES of published
numbers.}

\begin{table}[h]
\centering
\small
\begin{tabular}{lcccc}
\toprule
Task & Metric & \cite{zhang2024card} & Ours & $\Delta$ \\
\midrule
Line & EM & 52.56 & \todo{} & \todo{} \\
     & ES & 72.26 & \todo{} & \todo{} \\
API  & EM & 40.38 & \todo{} & \todo{} \\
     & ES & 67.01 & \todo{} & \todo{} \\
\bottomrule
\end{tabular}
\caption{CARD-RG1 with CodeLlama-7B on RepoEval. Our reimplementation vs the
published numbers.}
\label{tab:card_reproduction}
\end{table}

\subsection{RQ1: cascade vs. baselines}

\todo{Headline table: 6 configs $\times$ 3 datasets. For each cell:
EM, ES, IdF1, \% retrieval, mean latency.}

\begin{table*}[h]
\centering
\small
\begin{tabular}{lcccccc}
\toprule
Config & EM & ES & IdF1 & RSP & Hall.\% & Retr.\% \\
\midrule
C1 No-retrieve     & \todo{} & \todo{} & \todo{} & \todo{} & \todo{} & 0.0 \\
C2 Always-retrieve & \todo{} & \todo{} & \todo{} & \todo{} & \todo{} & 100.0 \\
C3 CARD            & \todo{} & \todo{} & \todo{} & \todo{} & \todo{} & \todo{} \\
\textbf{C4 Cascade (ours)} & \todo{} & \todo{} & \todo{} & \todo{} & \todo{} & \todo{} \\
C5 Static-only     & \todo{} & \todo{} & \todo{} & \todo{} & \todo{} & \todo{} \\
C6 Oracle          & \todo{} & \todo{} & \todo{} & \todo{} & \todo{} & \todo{} \\
\bottomrule
\end{tabular}
\caption{Headline results on CrossCodeEval-Python. RSP =
repository-symbol precision; Hall.\% = per-instance hallucination rate.}
\label{tab:rq1_cce}
\end{table*}

\todo{Narrative: H1 is supported if C4 matches C2 on ES within X and
beats C1 strictly. Report the cascade's retrieval savings vs C2 and
hallucination reduction vs C3. Discuss when the cascade ties C2
(probably: low ratio of cross-file dependencies) and when it beats C2
(high cross-file ratio, smaller models).}

\subsection{RQ2: hallucination reduction from static analysis}

We compare hallucination rate between vanilla CARD (C3) and the cascade
(C4) on the same 2{,}665 CrossCodeEval-Python instances, using McNemar's
exact test on the per-instance hallucination flag.

\todo{Report b, c, p-value. Hypothesis H2 is supported if c $\gg$ b
(cascade fixes more hallucinations than it introduces) and p $<$ 0.05.}

\subsection{RQ3: CARD vs. baselines}

\todo{Reproducing CARD's published findings: matches always-retrieve
within $\pm$ 1 ES while saving 20--46\% of retrievals.}

\subsection{Ablations}
\label{sec:ablations}

\paragraph{A1: per-signal contribution.} The four static signals
(\textbf{unresolved}, \textbf{signature}, \textbf{import},
\textbf{crossfile}) are independent flags; the headline cascade enables
the first three and disables the fourth. To measure each signal's
marginal contribution we run the cascade with one signal at a time,
report retrieval rate, ES, identifier-F1, and hallucination rate, and
report the leave-one-out drop relative to the full cascade.

\begin{table}[h]
\centering
\small
\begin{tabular}{lccccc}
\toprule
Variant & EM & ES & IdF1 & Hall.\% & Retr.\% \\
\midrule
unresolved only            & \todo{} & \todo{} & \todo{} & \todo{} & \todo{} \\
signature only             & \todo{} & \todo{} & \todo{} & \todo{} & \todo{} \\
import only                & \todo{} & \todo{} & \todo{} & \todo{} & \todo{} \\
crossfile only             & \todo{} & \todo{} & \todo{} & \todo{} & \todo{} \\
\midrule
unresolved + signature     & \todo{} & \todo{} & \todo{} & \todo{} & \todo{} \\
unresolved + import        & \todo{} & \todo{} & \todo{} & \todo{} & \todo{} \\
\textbf{full cascade (ours)} & \todo{} & \todo{} & \todo{} & \todo{} & \todo{} \\
full + crossfile           & \todo{} & \todo{} & \todo{} & \todo{} & \todo{} \\
\bottomrule
\end{tabular}
\caption{A1 ablation on CrossCodeEval-Python. The full cascade uses
unresolved + signature + import.}
\label{tab:a1_per_signal}
\end{table}

\paragraph{A2: $\TRAG$ sweep.} $\TRAG \in \{0.7, 0.8, 0.9\}$ on the
full cascade.

\paragraph{A3: Retrieval top-k.} top-k $\in \{5, 10, 20\}$ on the
full cascade.

\section{Analysis}
\label{sec:analysis}

\subsection{Per-trigger-reason breakdown}

The cascade now reports one of five trigger reasons per instance:
\texttt{none}, \texttt{card}, \texttt{static\_unresolved},
\texttt{static\_signature}, or \texttt{static\_import}
(plus \texttt{static\_crossfile} when the ablation flag is enabled).
For each reason we report the count, fraction of the test set, mean ES,
and per-instance hallucination rate of the resulting prediction. The
three \texttt{static\_*} buckets are the cleanest evidence for our
contribution: in each, CARD said ``skip'' yet the prediction's AST
exhibited one of the three structural hallucination signatures, and the
cascade's extra retrieval was meant to fix it.

\begin{table}[h]
\centering
\small
\begin{tabular}{lcccc}
\toprule
Trigger & Count & \% & ES & Hall.\% \\
\midrule
none                  & \todo{} & \todo{} & \todo{} & \todo{} \\
card                  & \todo{} & \todo{} & \todo{} & \todo{} \\
static\_unresolved    & \todo{} & \todo{} & \todo{} & \todo{} \\
static\_signature     & \todo{} & \todo{} & \todo{} & \todo{} \\
static\_import        & \todo{} & \todo{} & \todo{} & \todo{} \\
\bottomrule
\end{tabular}
\caption{C4 cascade trigger-reason breakdown on CrossCodeEval-Python.}
\label{tab:trigger_breakdown}
\end{table}

\subsection{Disagreement analysis: where CARD and cascade differ}

\todo{Focus on the ``card\_no\_cascade\_yes'' bucket: instances where
CARD's gate said skip but static analysis added a retrieval. Report
how often this bucket's ES improves under cascade and characterize
qualitatively what kinds of names triggered the static gate.}

\subsection{Threshold sweep}

\todo{Figure: ES, hallucination rate, retrieval \% as a function of
$\TRAG$. Reproduced by scripts/06\_analysis.py.}

\begin{figure}[h]
\centering
\includegraphics[width=\columnwidth]{figures/t_rag_sweep.pdf}
\caption{$\TRAG$ sweep on CrossCodeEval-Python with the CARD ŝ values from
the C3 run, paired with C1/C2 predictions. The cascade picks the operating
point that maximises ES at the lowest retrieval cost.}
\label{fig:trag_sweep}
\end{figure}

\subsection{Qualitative examples}

\todo{2--3 concrete instances pulled from the JSONL: one where the cascade's
static gate caught a hallucination CARD missed, one where the cascade
and CARD agreed, and one where the cascade triggered a retrieval that
didn't help (failure case). Show $\hat{y}_0$, the cascade's final
prediction, and the ground truth.}

\section{Discussion}

\paragraph{Unresolved.} The unresolved bucket is an over-approximation
of ``hallucinated identifier'': any name absent from our repo-wide
symbol table fires, but Python's \texttt{importlib}, dynamic attribute
access, and runtime name binding can produce legitimate code whose
identifiers we don't see statically. We mitigate this by populating
visible names from imports (including aliased and \texttt{from X import
Y} forms), function parameters, comprehension and walrus targets, with
and except aliases; the failures we expect to see in qualitative
analysis are runtime-only bindings.

\paragraph{Signature.} The signature gate is precision-biased by
construction: it abstains on method calls with untyped receivers
(\texttt{self.x}, \texttt{obj.x}), on multi-step attribute chains,
on calls to externally-defined functions, and on callees decorated with
signature-changing wrappers we don't recognise. The cost is reduced
recall --- a hallucinated method call on an in-file object will not
fire signature even though Tier 1 receiver-type inference would catch
it; we sketch that extension as future work. Conversely, the precision
gain is concrete: every fire is backed by a uniquely-resolved repo
signature.

\paragraph{Import.} The wrong-import gate handles re-exports via
\texttt{\_\_init\_\_.py} imports and \texttt{\_\_all\_\_} declarations,
but star imports and conditional or runtime-dynamic re-exports defeat
it. We abstain on these rather than risk false fires.

\paragraph{Cross-file.} Even with cross-file disabled, the
\texttt{cross\_file\_identifiers} field is still recorded per instance
so the Phase 7 disagreement analysis can mine cases where the model
recovered a cross-file name and the cascade was right not to fire.

\paragraph{Language portability.} Tree-sitter has parsers for most
languages, so the unresolved and signature stages are portable in
principle. The import stage is Python-specific in its handling of
\texttt{\_\_init\_\_.py} semantics; a Java or TypeScript port would
need analogous module-resolution logic.

\section{Conclusion}

We added a static-analysis cascade stage to CARD's adaptive-retrieval
framework. The new stage targets three orthogonal blind spots of CARD's
uncertainty signal --- \textbf{unresolved identifiers},
\textbf{signature mismatches}, and \textbf{wrong-import attributions} ---
and triggers retrieval on instances CARD would otherwise skip. Across
CrossCodeEval-Python and RepoEval, the cascade preserves CARD's
retrieval savings while reducing per-instance hallucination rate by
\todo{XX\%}. Prediction-side static analysis is a cheap, training-free,
language-portable complement to logit-based uncertainty: it costs one
tree-sitter parse per skipped instance and surfaces structural evidence
of hallucinations that logit aggregation provably cannot see.

\bibliographystyle{plain}
\bibliography{refs}

\end{document}
